\section{Experimental Design}
\label{sec:design}

The study has two parts. The \emph{main benchmark} (Sections~\ref{sec:design-suite}--\ref{sec:design-stats}) compares
Kumo Tabular with seven classical models on 21 public datasets, using the dataset as the unit of inference. The
\emph{case studies} (Section~\ref{sec:design-case}) are a re-implementation, with additional analyses, of the
author's earlier two-dataset study, and are reported as secondary, dataset-specific evidence. All case-study numbers
come from this re-implementation; the outputs of the earlier experiments are not available, and we make no claim that
they are reproduced here. The dataset selection criteria, splits, models and computed metrics
were fixed before the benchmark was run. The analysis plan---primary metrics, statistical tests and sensitivity
analyses---was written down and committed to a local version-control repository after a pipeline check on one dataset
(\texttt{wdbc}), whose complete results for all eight models we had seen, and before results for any other dataset
existed (supplementary material). The plan is therefore not a formal preregistration: it has no external timestamp,
and it was written with knowledge of one of the 21 datasets. One planned sensitivity analysis excludes \texttt{wdbc}.

\subsection{Datasets}
\label{sec:design-suite}
We took all datasets from OpenML-CC18 \citep{bischl2021openml} and OpenML-CTR23 \citep{fischer2023ctr23} that have
150 to 2{,}000 rows, at most 60 features, only numerical features, no missing values and, for classification, at most
ten classes (the limit of the Kumo Tabular classification head). The filters selected 15 classification and 7
regression datasets. One of them, \texttt{forest\_fires}, contained two text columns although its OpenML metadata
listed none; it was excluded under the same criterion when loading revealed this, before any of its results existed.
This leaves 15 classification and 6 regression datasets (Table~\ref{tab:datasets}). The restrictions define the regime we study---small, fully numerical tables---and
avoid preprocessing choices (encoding, imputation, subsampling) that would confound the comparison. They also limit the
scope of our conclusions (Section~\ref{sec:limitations}). One dataset, \texttt{wdbc}, is the Breast Cancer Wisconsin
data also used in the case studies.

\begin{table}[t]
\centering
\small
\caption{Datasets of the OpenML benchmark, selected from OpenML-CC18 and OpenML-CTR23 by fixed criteria (150--2{,}000 rows, at most 60 features, numerical features only, no missing values, at most 10 classes). \texttt{forest\_fires} (OpenML 44962) met the criteria according to its OpenML metadata but contains two text columns and was excluded.}
\label{tab:datasets}
\smallskip
\begin{adjustbox}{max width=\textwidth}
\begin{tabular}{rlrrr|rlrr}
\toprule
\multicolumn{5}{c|}{Classification} & \multicolumn{4}{c}{Regression} \\
ID & Name & Rows & Feat. & Cls. & ID & Name & Rows & Feat. \\
\midrule
11 & balance-scale & 625 & 4 & 3 & 44957 & airfoil\_self\_noise & 1503 & 5 \\
18 & mfeat-morphological & 2000 & 6 & 10 & 44959 & concrete\_compressive\_strength & 1030 & 8 \\
22 & mfeat-zernike & 2000 & 47 & 10 & 44960 & energy\_efficiency & 768 & 8 \\
37 & diabetes & 768 & 8 & 2 & 44970 & QSAR\_fish\_toxicity & 908 & 6 \\
54 & vehicle & 846 & 18 & 4 & 44972 & red\_wine & 1599 & 11 \\
1049 & pc4 & 1458 & 37 & 2 & 44994 & cars & 804 & 17 \\
1050 & pc3 & 1563 & 37 & 2 & & & & \\
1063 & kc2 & 522 & 21 & 2 & & & & \\
1068 & pc1 & 1109 & 21 & 2 & & & & \\
1462 & banknote-authentication & 1372 & 4 & 2 & & & & \\
1464 & blood-transfusion-service-center & 748 & 4 & 2 & & & & \\
1494 & qsar-biodeg & 1055 & 41 & 2 & & & & \\
1510 & wdbc & 569 & 30 & 2 & & & & \\
40982 & steel-plates-fault & 1941 & 27 & 7 & & & & \\
40994 & climate-model-simulation-crashes & 540 & 20 & 2 & & & & \\
\bottomrule
\end{tabular}
\end{adjustbox}
\end{table}


\subsection{Splits and seeds}
We use the official OpenML task splits: 10-fold cross-validation, stratified for classification, first repetition for
every task. The fold index (0--9) seeds every model and the random number generator of Kumo Tabular, whose predictions
are deterministic given the seed.

\subsection{Models}
All models receive the same folds and no hyperparameter tuning; the comparison is between out-of-the-box methods.
\begin{itemize}
\item \textbf{Kumo Tabular} \texttt{small}, 8 ensemble members, training fold as context, test fold as query, float32
inference (float16 autocast is disabled on GPUs older than compute capability 7.0).
\item \textbf{Logistic regression} on standardized inputs (\texttt{max\_iter}=5000) for classification;
\textbf{ordinary least squares} for regression.
\item \textbf{Random forest} and \textbf{extra trees}, 500 trees each.
\item \textbf{Histogram gradient boosting} (scikit-learn), \textbf{XGBoost}, \textbf{LightGBM} and \textbf{CatBoost},
library defaults.
\end{itemize}

\subsection{Metrics}
The primary metrics, fixed in advance, are the log loss for classification, a proper scoring rule for probabilistic
predictions \citep{gneiting2007scoring}, and the root mean squared error (RMSE) for regression. Secondary metrics are
accuracy, macro-averaged F1 and one-vs-rest macro-averaged ROC-AUC \citep{fawcett2006roc} for classification, and mean absolute error and
$R^2$ for regression. Each metric is computed per fold and averaged over the ten folds of a dataset. For models that do
not output a probability for a class absent from a training fold, that probability is zero.

\subsection{Statistical analysis}
\label{sec:design-stats}
Folds of one dataset share most of their training rows and are not independent, so we use the per-dataset mean as the
unit of analysis and treat datasets as independent replicates \citep{demsar2006comparisons}. For each task and metric
we apply a Friedman test \citep{friedman1937ranks} to the ranks of all eight models, and compare Kumo Tabular with each
baseline by a two-sided Wilcoxon signed-rank test on the per-dataset means, with Holm's correction \citep{holm1979}
across the seven baselines. Only the primary metrics support confirmatory statements. We report wins, ties and losses
and the median relative advantage as effect sizes. With six regression datasets the smallest attainable two-sided
Wilcoxon $p$-value is $2/2^6\approx0.031$, which exceeds $0.05$ after Holm's correction for seven comparisons: the
regression results are descriptive by design. The plan includes two sensitivity analyses: excluding \texttt{wdbc},
which overlaps with the case studies (and is the dataset seen before the plan was written), and averaging related datasets into single units (the NASA software-defect
datasets \texttt{pc1}, \texttt{pc3}, \texttt{pc4}, \texttt{kc2}; the two \texttt{mfeat} datasets) to reduce
dependence between datasets.

\subsection{Runtime}
We record wall-clock fit and prediction time per fold. Kumo Tabular has no fitting stage; all of its cost is
prediction. It ran on an NVIDIA GeForce GTX 980M (8\,GB) and the baselines on an Intel Core i7-6700HQ CPU, with one
untimed warm-up call per dataset for Kumo Tabular. Runtimes therefore compare complete pipelines on the available
hardware, not algorithms, and are reported descriptively.

\subsection{Case studies}
\label{sec:design-case}
The case studies use the scikit-learn copies of Breast Cancer Wisconsin (569 rows, 30 features, malignant as positive
class) and Diabetes (442 rows, 10 features, original units) \citep{street1993nuclear,efron2004lars}, five 80/20 splits
(seeds 42--46, stratified for classification), the first three baselines above, threshold metrics and RMSE, a
training-size study (25 to all training rows; seeds 42--44), the coverage of Kumo Tabular's nominal 90\% quantile
interval, and a within-dataset paired analysis with the Nadeau--Bengio corrected resampled $t$-test
\citep{nadeau2003inference} and Holm's correction. Because the splits of one dataset are dependent and only five, this
analysis is descriptive.

\subsection{Leakage controls and verification}
Kumo Tabular receives the training fold with labels and the test fold without its target column; unit tests check that
no test-row label reaches the model. OpenML splits are checked to be disjoint within each fold and to cover every row
exactly once. Preprocessing (standardization for logistic regression; Kumo Tabular's internal recipe) is fitted on the
training fold only. All case-study results were re-run with bit-identical metrics, and the case-study notebooks
\citep{anonymous2026supplement} were executed in a separate cloud environment (NVIDIA T4, PyTorch 2.11, float16 inference), where results for Kumo Tabular,
logistic and linear regression and random forests matched at the reported precision while histogram gradient boosting
differed slightly, presumably because of a different scikit-learn version.
